\chapter{Inherent and Passive Safety in Modern Designs}
\label{ch:inherent}

The lesson of Chernobyl---that safety must live in the physics, not in perfect
operation---has shaped reactor design ever since. This chapter surveys how modern
reactors build in negative feedback and passive heat removal so that thermal
stability and decay-heat safety are guaranteed by nature rather than by
intervention.

\section{Inherent vs.\ engineered safety}

\emph{Engineered} safety relies on active systems---pumps, valves, diesel
generators, operator actions---to respond to upsets. \emph{Inherent} safety relies
on the physical properties of the reactor: feedback coefficients, material
choices, and geometry that make dangerous states physically self-correcting or
unreachable. \emph{Passive} safety relies on natural forces---gravity, natural
circulation, conduction, stored coolant---that need no power or operator action.
The modern ideal combines all three with defence in depth, but assigns the
foundational role to inherent and passive features, precisely because Chernobyl,
TMI, and Fukushima showed that active systems and operators can fail.

\section{Designing in negative feedback}

The first principle is a \textbf{negative power coefficient across all operating
states}. This is achieved through:
\begin{itemize}[nosep]
\item \textbf{Under-moderation} in light-water reactors, guaranteeing negative
moderator and void coefficients (the opposite of the RBMK).
\item Strong \textbf{Doppler feedback} from fertile resonance absorbers
($^{238}$U, $^{240}$Pu), giving a prompt negative coefficient that limits
excursions.
\item Avoidance of positive void pathways: no graphite-decoupled-from-coolant
geometry, and shutdown systems that insert only negative reactivity.
\end{itemize}
Licensing in most countries requires demonstrating a negative power coefficient and
analysing reactivity-initiated accidents to show self-limitation. The RBMK's
positive void coefficient would not be licensable under such rules.

\section{Passive decay-heat removal}

Because decay heat (Chapter~\ref{ch:decay}) is irreducible, modern designs provide
\emph{passive} paths to remove it without power. Examples include gravity-driven
core flooding, natural-circulation cooling loops, large elevated water pools that
boil off over days, and---in some advanced designs---heat conduction to the ground
or air. The goal is a reactor that can survive a complete station blackout (as at
Fukushima) for an extended ``grace period'' with no operator action and no
electrical power. The AP1000, for instance, is designed to remove decay heat
passively for about three days using stored water and natural circulation.

\section{Generation III+ light-water reactors}

Evolutionary designs (AP1000, EPR, and others) retain the proven negative-feedback
physics of the PWR and add passive safety systems, larger water inventories,
core-catchers to contain molten fuel, and digital protection. They are, in
stability terms, conservative: their inherent feedback is strongly negative and
their innovations are mostly in passive heat removal and containment. They embody
the direct engineering response to the failure modes of TMI, Chernobyl, and
Fukushima.

\section{Generation IV and advanced reactors}

More radical designs pursue inherent safety through different physics:
\begin{itemize}[nosep]
\item \textbf{High-temperature gas reactors} (HTGR) use TRISO fuel particles that
retain fission products to very high temperatures and a core whose large negative
temperature coefficient and low power density allow it to shed decay heat by
conduction and radiation alone---it cannot melt even with no cooling.
\item \textbf{Molten-salt reactors} carry fuel dissolved in a liquid salt with a
strongly negative temperature/expansion coefficient and a passive
``freeze-plug'' drain that empties the core into safe geometry on loss of power.
\item \textbf{Sodium-cooled fast reactors} can be designed with strong negative
expansion feedback; the EBR-II tests in 1986---the same year as Chernobyl---famously
demonstrated survival of a complete loss of flow \emph{without scram} purely
through inherent feedback.
\item \textbf{Lead-cooled and heat-pipe microreactors} pursue similar
walk-away-safe goals with passive heat transport.
\end{itemize}
The unifying theme is the same principle this book has developed: make the
dominant feedback strongly negative and provide a passive path for decay heat, so
that the reactor is safe \emph{by physics}.

\section{The EBR-II demonstration: the anti-Chernobyl}

It is worth pausing on a remarkable coincidence of timing. On 3~April 1986---just
three weeks before Chernobyl---engineers at the Experimental Breeder Reactor~II
(EBR-II) in Idaho deliberately subjected the reactor to a complete
\emph{loss-of-flow-without-scram} test at full power: they cut the coolant pumps
and disabled the automatic shutdown. Rather than melting, the reactor's strong
\emph{negative} reactivity feedback (from thermal expansion of the metallic fuel
and structure) reduced the power as temperatures rose, and the reactor settled to a
safe, low-power state on its own. The same week, two reactors on opposite sides of
the world were pushed beyond their protection systems; one destroyed itself and a
region, the other quietly proved that inherent negative feedback is the surest
safety system of all. The contrast is the thesis of this book in a single image.

\keypoint{Modern reactor safety is founded on inherent negative feedback (a
negative power coefficient guaranteed across all states) and passive decay-heat
removal (natural forces that need no power or operators). EBR-II's 1986
loss-of-flow-without-scram test is the positive mirror image of Chernobyl:
inherent negative feedback turning a would-be accident into a non-event.}



\section{The quasi-static reactivity balance and integral safety parameters}
\label{sec:qsrb}

The qualitative claim that inherent feedback ``shuts the reactor down'' can be made
quantitative through the \emph{quasi-static reactivity balance} (QSRB) of Wade and
Fujita, the analytical backbone of metal-fuelled fast-reactor passive safety. For
transients slow compared with the prompt and delayed-neutron timescales, the
reactor passes through a sequence of near-critical states in which the net
reactivity is approximately zero. Writing $\tilde P=P/P_0$ for normalised power and
$\tilde F=F/F_0$ for normalised flow, the asymptotic balance reads
\begin{equation}
0=\delta\rho_{\rm ext}+(\tilde P-1)\,A+\Big(\tfrac{\tilde P}{\tilde F}-1\Big)B
   +\delta T_{\rm in}\,C,
\label{eq:qsrb}
\end{equation}
where the \emph{integral reactivity parameters} are
\begin{equation}
A=P_0\frac{\partial\rho}{\partial P}\ (\text{power coefficient}\times P_0),\quad
B=\text{power/flow coefficient},\quad
C=\frac{\partial\rho}{\partial T_{\rm in}}\ (\text{inlet-temperature coefficient}),
\end{equation}
with $A<0$ (Doppler $+$ axial fuel expansion) and $B<0$ (coolant-density/radial
expansion) for an inherently safe core. Equation \eqref{eq:qsrb} compresses all the
feedbacks of Chapter~\ref{ch:feedback} into three measurable integral numbers.

\begin{proposition}[Asymptotic power of an unprotected loss of flow]
\label{prop:ulof}
For an unprotected loss-of-flow transient (ULOF: $\delta\rho_{\rm ext}=0$,
$\delta T_{\rm in}=0$, flow decaying to $\tilde F<1$) the asymptotic power is
\begin{equation}
\tilde P_\infty=\frac{A+B}{A+B/\tilde F},
\end{equation}
which satisfies $0<\tilde P_\infty<1$ whenever $A,B<0$; the core settles to a
reduced power without scram. The asymptotic coolant temperature rise across the
core obeys $\Delta T_c^\infty/\Delta T_c^0=\tilde P_\infty/\tilde F$.
\end{proposition}
\begin{proof}
Set $\delta\rho_{\rm ext}=\delta T_{\rm in}=0$ in \eqref{eq:qsrb}:
$(\tilde P-1)A+(\tilde P/\tilde F-1)B=0$, i.e.\ $\tilde P(A+B/\tilde F)=A+B$, giving
$\tilde P_\infty$. With $A,B<0$ and $0<\tilde F<1$, both numerator and denominator
are negative and the denominator is larger in magnitude, so
$0<\tilde P_\infty<1$. The temperature ratio follows from
$\Delta T_c\propto P/F$.
\end{proof}

The dimensionless ratios $A/B$ and $C\,\Delta T_c/A$ are the modern \emph{passive
safety metrics}: a design with $A/B\approx1$ and small $|C\,\Delta T_c/A|$ rides
out the unprotected ULOF, ULOHS (loss of heat sink), and UTOP (transient
overpower) accidents with asymptotic temperatures below failure limits---the
analytic statement of the EBR-II demonstrations.

\section{Passive decay-heat removal: natural-circulation scaling}
\label{sec:natcirc}

Passive cooling replaces pumps with buoyancy. In a closed loop of thermal-centre
elevation difference $\Delta z$, the buoyancy head balances friction:
\begin{equation}
\underbrace{\rho_0\,\beta_{\rm th}\,g\,\Delta z\,\Delta T_{\rm loop}}_{\text{buoyancy}}
=\underbrace{R_f\,\dot m^2}_{\text{friction}},\qquad
R_f=\frac{1}{2\rho_0 A_{\rm fl}^2}\Big(\frac{fL}{D}+\textstyle\sum K\Big),
\end{equation}
with $\beta_{\rm th}$ the thermal expansion coefficient and
$\Delta T_{\rm loop}=Q/(\dot m c_p)$ the loop temperature rise carrying power $Q$.

\begin{proposition}[One-third power law]
\label{prop:natcirc}
Eliminating $\Delta T_{\rm loop}$ gives the natural-circulation mass flow and
temperature rise
\begin{equation}
\dot m=\Big(\frac{\rho_0\beta_{\rm th}g\Delta z}{c_p R_f}\Big)^{1/3}Q^{1/3},
\qquad
\Delta T_{\rm loop}=\frac{Q}{c_p\dot m}\propto Q^{2/3}.
\end{equation}
Hence the coolant temperature rise grows only as $Q^{2/3}$: as decay heat falls
after shutdown, natural circulation self-adjusts and the temperature rise decays
sublinearly, the quantitative basis of multi-day passive grace periods.
\end{proposition}
\begin{proof}
Substitute $\Delta T_{\rm loop}=Q/(\dot m c_p)$ into the head--friction balance:
$\rho_0\beta_{\rm th}g\Delta z\,Q/(\dot m c_p)=R_f\dot m^2$, so
$\dot m^3=\rho_0\beta_{\rm th}g\Delta z\,Q/(c_p R_f)$, giving the stated
$\dot m\propto Q^{1/3}$ and therefore $\Delta T_{\rm loop}=Q/(c_p\dot m)\propto
Q/Q^{1/3}=Q^{2/3}$.
\end{proof}

\section{Conduction-cooled cores and the radiative walk-away criterion}
\label{sec:walkaway}

In a high-temperature gas or microreactor core with no active cooling, decay heat
leaves by conduction to the vessel and thence by radiation to a passive heat sink.
Modelling the vessel as a grey radiator of area $A_s$ and emissivity
$\varepsilon$, energy balance at the quasi-steady decay power $P_d$ gives
\begin{equation}
\varepsilon\sigma A_s\big(T_{\rm eq}^4-T_{\rm amb}^4\big)=P_d
\ \Longrightarrow\
T_{\rm eq}=\Big(\frac{P_d}{\varepsilon\sigma A_s}+T_{\rm amb}^4\Big)^{1/4}.
\label{eq:radeq}
\end{equation}
Adding the internal conduction rise from a distributed source $q'''$ across a core
of conductivity $k$ and characteristic radius $R$ (cylindrical, parabolic profile)
$\Delta T_{\rm cond}=q''' R^2/(4k)=P_d/(8\pi k H)$ for height $H$, the peak fuel
temperature is $T_{\max}=T_{\rm eq}+\Delta T_{\rm cond}$.

\begin{definition}[Walk-away-safe core]
A core is \emph{walk-away safe} against a complete, unmitigated loss of cooling if
its peak passive temperature satisfies $T_{\max}<T_{\rm fail}$, where
$T_{\rm fail}$ is the fuel-failure temperature (e.g.\ $\sim1600\degC$ for TRISO
fuel). By \eqref{eq:radeq} this requires the power density to be low enough that
$T_{\rm eq}+P_d/(8\pi k H)<T_{\rm fail}$, a constraint that scales favourably as
core size shrinks ($A_s/V$ rises), explaining the passive safety of small modular
and micro-reactors.
\end{definition}

Because $T_{\rm eq}\propto P_d^{1/4}$ while the failure margin is fixed, the
walk-away criterion is fundamentally a \emph{power-density} limit: it caps the
admissible $P_d/A_s$, and through the decay-heat fraction (Chapter~\ref{ch:decay}),
the admissible operating power density---a recurring constraint in Generation-IV
design optimisation~\cite{waltar}.


\section*{Exercises}
\addcontentsline{toc}{section}{Exercises}
\begin{enumerate}[leftmargin=*,itemsep=2pt]
\item Distinguish inherent, passive, and engineered safety with an example of each.
\item Explain how under-moderation guarantees a negative void coefficient.
\item Describe one passive decay-heat removal strategy and the grace period it provides.
\item Summarise the EBR-II 1986 loss-of-flow-without-scram test and why it is the ``anti-Chernobyl''.
\item For an HTGR with TRISO fuel, explain why the core ``cannot melt'' even without active cooling.
\end{enumerate}
